% GENERATED FILE. Edit canonical lessons, then run npm run generate:books.
% canonical-source-hash: fnv1a64:99219847afca61cd
% canonical-lessons: KA-C452-sarige, KA-C452-dipo, KA-C452-saraku, KA-C452-sagisu, KA-C452-koriyar

\chapter{Moving Goods}
\label{ch:ka-moving-goods}
\hlchaptermodality{452}

\begin{chapteropening}
\textbf{By the end of this chapter:} \emph{I can say transport, a depot, goods, to transport and a courier.}

Complete the last lesson of chapter 452: I can say transport, a depot, goods, to transport and a courier.
\end{chapteropening}

\section[\texorpdfstring{\kn{ಸಾರಿಗೆ} (sārige)}{ಸಾರಿಗೆ (sārige)}]{\kn{ಸಾರಿಗೆ} (sārige) — transport; public transport}

\label{lesson:KA-C452-sarige}



\begin{quote}
Before the new one: say the Kannada for a motorbike, then the Kannada for a bullock cart.
\end{quote}

\subsection*{You'll want to know: \kn{ಸಾರಿಗೆ}}
\textbf{\kn{ಸಾರಿಗೆ}} — \emph{sārige} — \textquotedblleft{}transport; public transport\textquotedblright{}.

\begin{quote}
\kn{ಈ} \kn{ಊರಿಗೆ} \kn{ಸರ್ಕಾರಿ} \kn{ಸಾರಿಗೆ} \kn{ಬಸ್} \kn{ದಿನಕ್ಕೆ} \kn{ಎರಡು} \kn{ಸಲ} \kn{ಬರುತ್ತದೆ}.
\end{quote}

\begin{quote}
\emph{ī ūrige sarkāri sārige bas dinakke eraḍu sala baruttade.}
\end{quote}

\begin{quote}
A government transport bus comes to this town twice a day.
\end{quote}

\kn{ಸಾರಿಗೆ} \kn{ಇಲಾಖೆ}, the transport department.

\begin{cousinweb}
\textbf{\kn{ಸಾರಿಗೆ}} (\emph{sārige}) is a Kannada word for transport whose origin is not settled.
\end{cousinweb}

\begin{grammarlens}[title={What you've built}]
One more word for getting around.
\end{grammarlens}

\subsection*{Your turn}
\begin{itemize}
\raggedright
  \item \emph{Say it:} \emph{sārige}
  \item \emph{Say it:} \emph{sārige}, once more
  \item \emph{Say it:} \emph{ī ūrige sarkāri sārige bas dinakke eraḍu sala baruttade.} — A government transport bus comes to this town twice a day.
  \item \emph{Recall:} say the Kannada for a motorbike, then the Kannada for a bullock cart, then say \emph{sārige} again
\end{itemize}

\begin{tcolorbox}[breakable,colback=teal!4,colframe=teal!35!black,title={Before you move on}]
What does \kn{ಸಾರಿಗೆ} mean? (\textquotedblleft{}Transport; public transport\textquotedblright{}.) Say it once more.
\end{tcolorbox}
\section[\texorpdfstring{\kn{ಡಿಪೋ} (ḍipō)}{ಡಿಪೋ (ḍipō)}]{\kn{ಡಿಪೋ} (ḍipō) — a depot; a bus depot}

\label{lesson:KA-C452-dipo}



\begin{quote}
Before the new one: say the Kannada for a bullock cart, then the Kannada for transport.
\end{quote}

\subsection*{You'll want to know: \kn{ಡಿಪೋ}}
\textbf{\kn{ಡಿಪೋ}} — \emph{ḍipō} — \textquotedblleft{}a depot; a bus depot\textquotedblright{}.

\begin{quote}
\kn{ರಾತ್ರಿ} \kn{ಎಲ್ಲಾ} \kn{ಬಸ್ಸುಗಳು} \kn{ಡಿಪೋಗೆ} \kn{ಹಿಂತಿರುಗುತ್ತವೆ}.
\end{quote}

\begin{quote}
\emph{rātri ellā bassugaḷu ḍipōge hintiruguttave.}
\end{quote}

\begin{quote}
At night all the buses go back to the depot.
\end{quote}

Where buses are kept and mended.

\begin{cousinweb}
\textbf{\kn{ಡಿಪೋ}} (\emph{ḍipō}) is a loan from English \emph{depot}.
\end{cousinweb}

\begin{grammarlens}[title={What you've built}]
One more word for getting around.
\end{grammarlens}

\subsection*{Your turn}
\begin{itemize}
\raggedright
  \item \emph{Say it:} \emph{ḍipō}
  \item \emph{Say it:} \emph{ḍipō}, once more
  \item \emph{Say it:} \emph{rātri ellā bassugaḷu ḍipōge hintiruguttave.} — At night all the buses go back to the depot.
  \item \emph{Recall:} say the Kannada for a bullock cart, then the Kannada for transport, then say \emph{ḍipō} again
\end{itemize}

\begin{tcolorbox}[breakable,colback=teal!4,colframe=teal!35!black,title={Before you move on}]
What does \kn{ಡಿಪೋ} mean? (\textquotedblleft{}A depot; a bus depot\textquotedblright{}.) Say it once more.
\end{tcolorbox}
\section[\texorpdfstring{\kn{ಸರಕು} (saraku)}{ಸರಕು (saraku)}]{\kn{ಸರಕು} (saraku) — goods; cargo}

\label{lesson:KA-C452-saraku}



\begin{quote}
Before the new one: say the Kannada for transport, then the Kannada for a depot.
\end{quote}

\subsection*{You'll want to know: \kn{ಸರಕು}}
\textbf{\kn{ಸರಕು}} — \emph{saraku} — \textquotedblleft{}goods; cargo\textquotedblright{}.

\begin{quote}
\kn{ಲಾರಿಯಲ್ಲಿ} \kn{ಸರಕು} \kn{ತುಂಬಿದೆ}.
\end{quote}

\begin{quote}
\emph{lāriyalli saraku tumbide.}
\end{quote}

\begin{quote}
The lorry is full of goods.
\end{quote}

\kn{ಸರಕು} \kn{ರೈಲು}, a goods train.

\begin{cousinweb}
\textbf{\kn{ಸರಕು}} (\emph{saraku}) is native Dravidian: DEDR 2353 lists Kannada \emph{saraku}, \emph{sarku}, \textquotedblleft{}goods, things, merchandise, cargo\textquotedblright{}.
\end{cousinweb}

\begin{grammarlens}[title={What you've built}]
One more word for getting around.
\end{grammarlens}

\subsection*{Your turn}
\begin{itemize}
\raggedright
  \item \emph{Say it:} \emph{saraku}
  \item \emph{Say it:} \emph{saraku}, once more
  \item \emph{Say it:} \emph{lāriyalli saraku tumbide.} — The lorry is full of goods.
  \item \emph{Recall:} say the Kannada for transport, then the Kannada for a depot, then say \emph{saraku} again
\end{itemize}

\begin{tcolorbox}[breakable,colback=teal!4,colframe=teal!35!black,title={Before you move on}]
What does \kn{ಸರಕು} mean? (\textquotedblleft{}Goods; cargo\textquotedblright{}.) Say it once more.
\end{tcolorbox}
\section[\texorpdfstring{\kn{ಸಾಗಿಸು} (sāgisu)}{ಸಾಗಿಸು (sāgisu)}]{\kn{ಸಾಗಿಸು} (sāgisu) — to transport; to carry off}

\label{lesson:KA-C452-sagisu}



\begin{quote}
Before the new one: say the Kannada for a depot, then the Kannada for goods.
\end{quote}

\subsection*{You'll want to know: \kn{ಸಾಗಿಸು}}
\textbf{\kn{ಸಾಗಿಸು}} — \emph{sāgisu} — \textquotedblleft{}to transport; to carry off\textquotedblright{}.

\begin{quote}
\kn{ರೈತರು} \kn{ತಮ್ಮ} \kn{ಭತ್ತವನ್ನು} \kn{ಲಾರಿಯಲ್ಲಿ} \kn{ಸಾಗಿಸಿದರು}.
\end{quote}

\begin{quote}
\emph{raitaru tamma bhattavannu lāriyalli sāgisidaru.}
\end{quote}

\begin{quote}
The farmers carried their paddy away by lorry.
\end{quote}

Past (I …): \kn{ಸಾಗಿಸಿದೆ}, \emph{sāgiside}.

\kn{ಸಾಗು}, to move along, with -\kn{ಇಸು}: to make go.

\begin{cousinweb}
\textbf{\kn{ಸಾಗಿಸು}} (\emph{sāgisu}) is native Dravidian: the causative of \emph{sāgu}, \textquotedblleft{}to go, to move forward, to proceed\textquotedblright{} (DEDR 2430).
\end{cousinweb}

\begin{grammarlens}[title={What you've built}]
One more word for getting around.
\end{grammarlens}

\subsection*{Your turn}
\begin{itemize}
\raggedright
  \item \emph{Say it:} \emph{sāgisu}
  \item \emph{Say it:} \emph{sāgisu}, once more
  \item \emph{Say it:} \emph{raitaru tamma bhattavannu lāriyalli sāgisidaru.} — The farmers carried their paddy away by lorry.
  \item \emph{Recall:} say the Kannada for a depot, then the Kannada for goods, then say \emph{sāgisu} again
\end{itemize}

\begin{tcolorbox}[breakable,colback=teal!4,colframe=teal!35!black,title={Before you move on}]
What does \kn{ಸಾಗಿಸು} mean? (\textquotedblleft{}To transport; to carry off\textquotedblright{}.) Say it once more.
\end{tcolorbox}
\section[\texorpdfstring{\kn{ಕೊರಿಯರ್} (koriyar)}{ಕೊರಿಯರ್ (koriyar)}]{\kn{ಕೊರಿಯರ್} (koriyar) — a courier}

\label{lesson:KA-C452-koriyar}



\begin{quote}
Before the new one: say the Kannada for goods, then the Kannada for to transport.
\end{quote}

\subsection*{You'll want to know: \kn{ಕೊರಿಯರ್}}
\textbf{\kn{ಕೊರಿಯರ್}} — \emph{koriyar} — \textquotedblleft{}a courier\textquotedblright{}.

\begin{quote}
\kn{ನಾನು} \kn{ಊರಿಗೆ} \kn{ಒಂದು} \kn{ಪಾರ್ಸೆಲ್} \kn{ಕೊರಿಯರ್ನಲ್ಲಿ} \kn{ಕಳಿಸಿದೆ}.
\end{quote}

\begin{quote}
\emph{nānu ūrige ondu pārsel koriyarnalli kaḷiside.}
\end{quote}

\begin{quote}
I sent a parcel home by courier.
\end{quote}

\kn{ಕೊರಿಯರ್} \kn{ಮಾಡು}, to send by courier.

\begin{cousinweb}
\textbf{\kn{ಕೊರಿಯರ್}} (\emph{koriyar}) is a loan from English \emph{courier}.
\end{cousinweb}

\begin{grammarlens}[title={What you've built}]
One more word for getting around.
\end{grammarlens}

\subsection*{Your turn}
\begin{itemize}
\raggedright
  \item \emph{Say it:} \emph{koriyar}
  \item \emph{Say it:} \emph{koriyar}, once more
  \item \emph{Say it:} \emph{nānu ūrige ondu pārsel koriyarnalli kaḷiside.} — I sent a parcel home by courier.
  \item \emph{Recall:} say the Kannada for goods, then the Kannada for to transport, then say \emph{koriyar} again
\end{itemize}

\begin{tcolorbox}[breakable,colback=teal!4,colframe=teal!35!black,title={Before you move on}]
What does \kn{ಕೊರಿಯರ್} mean? (\textquotedblleft{}A courier\textquotedblright{}.) Say it once more.
\end{tcolorbox}
